\makeatletter
\def\input@path{{./}{../}{../../}{../../../}{../../../../}{preamble/}{../preamble/}{../../preamble/}{../../../preamble/}}
\makeatother

\input{preamble_exams}
\input{graphs/G11_Histograms.tex}

\graphicspath{{./}{../}{../../}{../../../}{../../../../}{preamble/}{../preamble/}{../../preamble/}{../../../preamble/}}

\setlength{\columnsep}{5mm}


\begin{document}
% HEADER BOX (3 COLUMNS)
\ExamHeader{11th}{C5, C6 and C7}{Continuous distributions and sample data}{1st}
\vspace{0.15cm}
% OBJECTIVE + CRITERIA
\LearningObjective
{Use continuous probability models and sample data to describe typical values,
variability, and interval probabilities, and interpret differences between
sample estimates and theoretical values.}
{C5: Calculates and interprets the mean, variance and standard deviation of continuous distributions using standard formulas and functions.\par
C6: Estimates interval probabilities using relative frequencies from generated samples.
 and compares them with theoretical probabilities.\par
C7: Calculates sample means, variances, and standard deviations for different sample sizes and data representations.}

% STUDENT LINE ---------------------------------
\vspace{0.1cm}
\StudentInfo
\vspace{-3mm}
% EXAM BODY
\begin{multicols}{2}
	\raggedcolumns
	
\SubsectionBox{Criteria assessment}\vspace{-0.2cm}
This exam assesses C5, C6, and C7 separately (6 points each; 18 total). Each task specifies its points. Partial credit requires correct procedures; answers without procedures receive zero. Missing names or late submissions incur a 1-point deduction per criterion.\textbf{Passing requires at least 4 points per criterion}. Use unrounded intermediate values and round final statistics to 3 decimals and percentage errors to 2.

\vspace{-9mm}
\begin{center}
	\begin{tabular}{|c|c|c|c|}
		\hline
		\textbf{Question} & \textbf{C5} & \textbf{C6} & \textbf{C7} \\ \hline
		1  & 2  & -- & -- \\ \hline
		2  & 2  & -- & -- \\ \hline
		3  & 2  & -- & -- \\ \hline
		4  & -- & 2  & -- \\ \hline
		5  & -- & 1  & -- \\ \hline
		6  & -- & 2  & -- \\ \hline
		7  & -- & 1  & -- \\ \hline
		8  & -- & -- & 2  \\ \hline
		9  & -- & -- & 1  \\ \hline
		10 & -- & -- & 1  \\ \hline
		11 & -- & -- & 2  \\ \hline
		\textbf{Total} & \textbf{6} & \textbf{6} & \textbf{6} \\ \hline
	\end{tabular}
\end{center}

\vspace{-5mm}
\SubsectionBox{1. Two custom-print studios \score{C5-\sml{2}}}
\phantomsection\label{sec:c5-1}
Two studios print custom posters, with each order taking between 3 and 15 hours
to complete. Both studios handle a wide variety of orders, but their production
processes differ.
Studio I specializes in large, detailed posters that require extensive
customization. As a result, longer completion times are more common, with
orders becoming progressively more likely to require additional production
time.
Studio D uses a more efficient printing process designed for quick orders.
Shorter completion times are more common, and orders become progressively
less likely to require additional production time.
Let $I$ and $D$ represent the respective completion times in hours.
Identify each distribution from the contextual information and calculate
its mean, variance, and standard deviation using the corresponding formulas.
\score{2}


\SubsectionBox{2. Bookings and daily takings \score{C5-\sml{2}}}
\phantomsection\label{sec:c5-2}

A camera repair shop's daily revenue varies due to many small, independent
factors, including customer demand and repair costs. Most days, revenue
is close to \$3200, while higher and lower amounts become increasingly
uncommon at similar rates. The standard deviation is \$200.
The shop receives online bookings independently at an average rate of
10 per hour. Bookings arrive at a steady average rate, and the time since
the last booking does not affect when the next arrives.
Identify the distributions for revenue and time between bookings and calculate their means,
variances, and standard deviations, expressing revenue in dollars and
waiting time in minutes. Interpret the mean revenue. \score{2}


\newpage

\SubsectionBox{3. Two parcel-labeling stations \score{C5-\sml{2}}}
\phantomsection\label{sec:c5-3}
A delivery company uses two parcel-labeling stations. At station A, an
automated machine labels parcels in 5 to 19 minutes. Due to differences
in parcel sizes and handling requirements, completion times vary
unpredictably, with no particular duration being more common than another.

At station B, staff label special parcels, also taking between 5 and
19 minutes. Most parcels are completed in around 10 minutes, while
shorter and longer processing times become progressively less common.
Very quick jobs are especially rare, and longer jobs occasionally
require additional handling.

Identify the distributions for the labeling times in both stations and calculate their respective means, variances, and standard deviations
using the corresponding formulas. \score{2}


\SubsectionBox{4. File upload times, grouped data \score{C6-\sml{2}}}
\phantomsection\label{sec:c6-1}
A cloud storage company monitors the time required for users to upload files
to its servers. To evaluate the performance of its system, the company
records the upload times, in seconds, for a sample of 24 files. The results
are summarized in the following table, where each class includes its lower
bound but excludes its upper bound.
\[
\begin{array}{c|rrrr}
	\text{Class}&[2,8)&[8,14)&[14,20)&[20,26)\\\hline
	f_i&5&9&7&3
\end{array}
\]
According to the company's theoretical model, the probability that a file
takes between 8 seconds and 20 seconds to upload is
0.75. The company wants to assess how closely the observed results agree
with this prediction. Calculate the empirical probability of this interval
and determine its absolute and percentage errors relative to the
theoretical probability. \score{2}


\SubsectionBox{5. Office electricity use, raw data \score{C6-\sml{1}}}
\phantomsection\label{sec:c6-2}
An office building monitors its daily electricity consumption to evaluate
how accurately its energy forecasting model predicts actual usage.
The building manager records electricity consumption over 20 working days,
measured in tens of kilowatt-hours. The observations are
\[
\begin{array}{cccccccccc}
	1.3&1.7&2.0&2.2&2.6&2.9&3.3&3.5&3.8&4.2\\
	1.6&1.9&2.3&2.7&3.1&3.4&3.6&3.9&4.4&4.6
\end{array}
\]
According to the building's theoretical energy model, the probability
that daily electricity consumption falls between 20 and 35 kWh,
including both limits, is 0.60. The manager wants to determine how
closely the observed consumption agrees with this prediction.
Calculate the empirical probability of this interval and determine
its absolute and percentage errors relative to the theoretical probability.\score{1}


\SubsectionBox{6. Museum entrance queues, histogram data \score{C6-\sml{2}}}
\phantomsection\label{sec:c6-3}
A museum monitors how long visitors wait in line before entering
to evaluate the accuracy of its queue management model.
During a busy weekend, staff record the waiting times of 40 visitors,
measured in minutes. The original observations are no longer available,
but their frequencies are preserved in the following histogram.
Each bar represents the number of visitors whose waiting times fall
within the corresponding interval, including the lower bound but
excluding the upper bound.
\begin{center}
	\HistogramGraph{(1,5),(5,9),(9,13),(13,17)}{8,15,11,6}
	{0.90\linewidth}{4cm}{blue!25}{4}{5}
\end{center}\vspace{-4mm}
According to the museum's theoretical model, the probability that
a visitor waits at least 5 minutes but less than 13 minutes is 0.70.
The museum wants to assess how closely the observed waiting times
agree with this prediction. Calculate the empirical probability
of this interval and determine its absolute and percentage errors
relative to the theoretical probability. \score{2}


\SubsectionBox{7. Equipment rental times, frequency data \score{C6-\sml{1}}}
\phantomsection\label{sec:c6-4}
An equipment rental company monitors how long customers use rented
equipment to evaluate the accuracy of its rental duration model.
The company records the rental times of 100 transactions, measured
in minutes. The results are summarized in the following frequency
table, where each value represents an observed rental duration
and its corresponding frequency.
\[
\begin{array}{c|rrrrrr}
	x_i\text{ (min)}&4&8&12&16&20&24\\\hline
	f_i&8&16&28&22&14&12
\end{array}
\]
According to the company's theoretical model, the probability that
a rental lasts at least 8 minutes but less than 20 minutes is 0.80.
The company wants to assess how closely the observed rental durations
agree with this prediction. Calculate the empirical probability
of this interval and determine its absolute and percentage errors
relative to the theoretical probability. \score{1}

\SubsectionBox{8. Luggage masses, histogram data \score{C7-\sml{2}}}
\phantomsection\label{sec:c7-1}
An airline monitors the weight of passengers' carry-on luggage to
evaluate the accuracy of its baggage weight model. During a routine
inspection, airport staff record the masses of 9 randomly selected
pieces of luggage, measured in kilograms. The individual measurements
are no longer available, but their frequencies are preserved in the
following histogram. Each bar represents the number of pieces of
luggage whose masses fall within the corresponding interval,
including the lower bound but excluding the upper bound.
\begin{center}
	\HistogramGraph{(4,7),(7,10),(10,13)}{3,4,2}
	{0.90\linewidth}{4.5cm}{orange!30}{3}{2}
\end{center}
According to the airline's theoretical model, luggage masses have
a mean of $\mu=9$ kg and a variance of $\sigma^2=7$ kg$^2$.
The airline wants to assess how closely the observed luggage
masses agree with these predictions. Using the grouped data,
estimate the sample mean, sample variance, and sample standard
deviation, then calculate the absolute and percentage errors
of each statistic relative to its corresponding theoretical value. \score{2}



\SubsectionBox{9. Chamber temperatures, raw data \score{C7-\sml{1}}}
\phantomsection\label{sec:c7-2}
A laboratory monitors the temperature inside a climate-controlled chamber
to evaluate the accuracy of its temperature regulation model. During a
routine inspection, technicians record five temperature readings, in
degrees Celsius: $14, 17, 18, 19, 22$.
According to the laboratory's theoretical model, the chamber temperature
has a mean of $\mu=19\,^\circ\mathrm{C}$ and a variance of
$\sigma^2=6\,^\circ\mathrm{C}^2$. The technicians want to assess how
closely the observed temperatures agree with these predictions.
Calculate the sample mean, sample variance, and sample standard
deviation, then determine the absolute and percentage errors of each
statistic relative to its corresponding theoretical value. \score{1}



\SubsectionBox{10. Parcels dispatched, frequency data \score{C7-\sml{1}}}
\phantomsection\label{sec:c7-3}
A delivery company monitors the number of parcels dispatched during
short collection periods to evaluate the accuracy of its logistics
forecasting model. During a routine assessment, the company records
the number of parcels dispatched across nine collection periods.
The results are summarized in the following frequency table,
where each value represents the number of parcels dispatched
and its corresponding frequency.
\[
\begin{array}{c|rrr}
	x_i&3&7&11\\\hline
	f_i&2&4&3
\end{array}
\]
According to the company's theoretical model, the number of parcels
dispatched per collection period has a mean of $\mu=7$ parcels
and a variance of $\sigma^2=9$ parcels$^2$. The company wants
to assess how closely the observed dispatch quantities agree
with these predictions. Calculate the sample mean, sample variance,
and sample standard deviation, then determine the absolute and
percentage errors of each statistic relative to its corresponding
theoretical value. \score{1}


\SubsectionBox{11. Walking times, grouped data \score{C7-\sml{2}}}
\phantomsection\label{sec:c7-4}
A university monitors how long students take to walk from the main
entrance to their classrooms to evaluate the accuracy of its campus
mobility model. During a routine assessment, staff record the walking
times of nine students, measured in minutes. The results are summarized
in the following frequency table, where each interval includes its
lower bound but excludes its upper bound.
\[
\begin{array}{c|rrr}
	\text{Time (min)}&[12,18)&[18,24)&[24,30)\\\hline
	f_i&2&3&4
\end{array}
\]
According to the university's theoretical model, walking times have
a mean of $\mu=23$ minutes and a variance of $\sigma^2=30$ minutes$^2$.
The university wants to assess how closely the observed walking times
agree with these predictions. Using the grouped data, estimate the
sample mean, sample variance, and sample standard deviation, then
calculate the absolute and percentage errors of each statistic
relative to its corresponding theoretical value. \score{2}


\end{multicols}
\end{document}
